\documentclass{article}
\usepackage{amsmath,amssymb,xurl}
\usepackage[hidelinks]{hyperref}
% Shared commands for notes in docs/paper-gaps/.

\DeclareUnicodeCharacter{2080}{\ifmmode{}_{0}\else\textsubscript{0}\fi}
\DeclareUnicodeCharacter{2081}{\ifmmode{}_{1}\else\textsubscript{1}\fi}
\DeclareUnicodeCharacter{2082}{\ifmmode{}_{2}\else\textsubscript{2}\fi}
\DeclareUnicodeCharacter{2099}{\ifmmode{}_{n}\else\textsubscript{n}\fi}

\newcommand{\Lean}{\textsc{Lean}}
\ExplSyntaxOn
\NewDocumentCommand{\leanid}{m}{
  \ifmmode
    \textsf{\detokenize{#1}}
  \else
    \group_begin:
    \urlstyle{sf}
    \tl_set:Nn \l_tmpa_tl {#1}
    \tl_replace_all:Nnn \l_tmpa_tl {\_} {_}
    \exp_args:NV \path \l_tmpa_tl
    \group_end:
  \fi
}
\ExplSyntaxOff
\providecommand{\tr}{\operatorname{tr}}
\newcommand{\tnleanIssueBase}{https://github.com/LionSR/TNLean/issues/}
\newcommand{\tnleanPrBase}{https://github.com/LionSR/TNLean/pull/}
\newcommand{\tnleanDocBase}{https://sirui-lu.com/TNLean/docs/}
\newcommand{\ghissue}[1]{\href{\tnleanIssueBase#1}{GitHub issue \##1}}
\newcommand{\ghpr}[1]{\href{\tnleanPrBase#1}{PR \##1}}
\newcommand{\doclink}[1]{\href{\tnleanDocBase#1.html}{\path{#1}}}

% Dirac notation and the identity operator, as in blueprint/src/macros/common.tex.
% \providecommand keeps note-local definitions authoritative.
\usepackage{dsfont}
\providecommand{\ket}[1]{|#1\rangle}
\providecommand{\bra}[1]{\langle #1|}
\providecommand{\braket}[2]{\langle #1 | #2 \rangle}
\providecommand{\ketbra}[2]{|#1\rangle\!\langle #2|}
\providecommand{\Id}{\mathds{1}}
\providecommand{\one}{\mathds{1}}
\providecommand{\C}{\mathbb{C}}
\providecommand{\MN}[1]{M_{#1}(\C)}
\providecommand{\MD}{M_D(\C)}

% Verdict marker: \gapnote{<kind>}{<status>}. Kind is the policy
% classification (clarification, local-correction, scope-restriction,
% unfaithful, false-source, open-gap); status is open, wip (elimination
% underway), resolved, or historical. Machine-read by the site index;
% prints a verdict line.
\newcommand{\gapnote}[2]{\par\noindent\textbf{Verdict:} #1 (#2).\par}

\title{Simply Connected Forbidden Regions and Torus Walk Detours}
\author{}
\date{}
\begin{document}
\maketitle
\gapnote{scope-restriction}{open}

Schuch, Cirac, and P\'erez-Garc\'ia,
\href{https://arxiv.org/abs/1001.3807v3}{arXiv:1001.3807v3},
Lemma 6.14, source \path{Papers/1001.3807/paper_v3.tex}, lines 2199--2214,
deforms a fixed-endpoint flux string. The exterior-collar argument in the
proof of Theorem 6.9, lines 1935--1990, gives geometric detours with
prescribed seam crossing numbers.

On a native torus with both periods at least three, let $R$ have simply
connected closed-cell realization. For every finite graph walk $p:a\to b$
with $a,b\notin R$, there exists a walk $q:a\to b$ in the induced complement
whose two integer seam crossing numbers agree with those of $p$. No exterior
walk or winding identity is supplied. The proof replaces each maximal
interior segment by a genuine collar detour and inducts on the remaining
length. The covering lift and collar connectivity are derived from simple
connectedness of the actual region.

The declaration is
\leanid{TNLean.PEPS.exists_complementWalk_same_winding_of_isSimplyConnected}.
This gives arbitrary-shape simply connected region avoidance for torus
site-graph walks. It does not yet identify these walks with the four-labelled
native dual paths, impose the physical crossed-bond incidence condition, or
prove vacuum reduced-density equality for open flux pairs. Those steps remain
under \ghissue{8703}. Arbitrary forbidden regions need not admit such a detour:
they may disconnect the endpoints or obstruct the requested winding class.
The full-lattice Hamiltonian work under \ghissue{8947} is unchanged.
\end{document}
